\documentclass[a4paper]{article}
\usepackage{exercise}
%um nur aufgaben zu zeigen
%\usepackage[noanswer]{exercise} 
\usepackage{../images/preamble}
\usepackage{rotating}
\usetikzlibrary{decorations.pathmorphing}
\usetikzlibrary{decorations.markings}
\usetikzlibrary{arrows}
\usetikzlibrary{shapes.geometric}
\newcommand{\midarrow}{\tikz \draw[-triangle 90] (0,0) -- +(.02,0);}
\usepackage{xcolor}

\usepackage{mathtools}


\newcommand{\re}{/tikz/circuitikz/bipoles/length=20pt, R}

%\usepackage[]{xwatermark}
%\newsavebox\mybox
%\savebox\mybox{\tikz[color=red,opacity=0.3]\node{\textsc{Entwurf}};}
%\newwatermark*[
%allpages,
%angle=45,
%scale=10,
%xpos=-4cm,
%ypos=4cm
%]{\usebox\mybox}
\pagestyle{fancy}
\fancyhead[L]{Aaron Wild}
\fancyhead[R]{\textit{Rekursives Widerstandsnetz}}


\begin{document}
%	\vspace*{-1cm}
%	\parbox{4cm}{\vspace{-0.2cm}\includegraphics[width=5cm]{../images/logo_scaled.pdf}}
%	\parbox{10.6cm}{\setstretch{2.0} \centering{ \huge \textsf{Neue Aufgaben?
%			}}\\ ~\\ \vspace*{-.5cm} }
%		\vspace{0.5cm}




\noindent

\input{../tasks/selfmade/initeloop}



\end{document}